% ============================================
% LUC-Contraprova
% Gerado de molambudos.md — build_miolo.py
% ============================================
\fragtitle{luc}{L}{LUC}{Contraprova}{A Contraprova}
\markboth{LUC-Contraprova --- A Contraprova}{}

\noindent \textit{\textquotedblleft{}Isso tem que ter uma explicação científica\textquotedblright{}}

\noindent \textbf{Abril de 2026 --- Belo Horizonte}

\vspace{0.5em}\noindent\rule{\textwidth}{0.5pt}\vspace{0.5em}

\noindent Três dias depois do sonho, Lúcia Mendes entrou no laboratório do Hospital das Clínicas da UFMG com uma pasta cheia de exames.

\noindent \textit{\textquotedblleft{}Preciso de uma bateria completa\textquotedblright{}\textit{, disse ao técnico. }\textquotedblleft{}Ressonância de crânio com contraste. Mapeamento retiniano. Eletroencefalograma de 24 horas. Hemograma completo. Dosagem hormonal. Perfil hepático. Teste de função tireoidiana. RPR. FAN. Toxoplasmose. Citomegalovírus. Teste rápido para HIV. Sífilis. Tudo.\textquotedblright{}}

\noindent O técnico ergueu uma sobrancelha. \textit{\textquotedblleft{}A senhora está com o quê?\textquotedblright{}}

\noindent \textit{\textquotedblleft{}Não sei. Quero descobrir.\textquotedblright{}}

\vspace{0.5em}\noindent\rule{\textwidth}{0.5pt}\vspace{0.5em}

\noindent Os resultados chegaram uma semana depois.

\noindent \textbf{Ressonância magnética:} Normal. Sem lesões, sem tumores, sem alterações de sinal no parênquima cerebral. Hipotálamo, amígdala, hipocampo --- todas as estruturas dentro dos padrões.

\noindent \textbf{Eletroencefalograma:} Atividade elétrica normal. Sem descargas epileptiformes. Sem padrões de atividade lenta focal. Sem evidência de atividade epileptogênica no lobo temporal.

\noindent \textbf{Mapeamento retiniano:\textit{ Epitélio pigmentar da retina normal. Disco óptico sem alterações. Pressão intraocular normal. }\textquotedblleft{}Seu olho direito está saudável\textquotedblright{}\textit{, disse o oftalmologista. }\textquotedblleft{}A descoloração não tem causa orgânica detectável.\textquotedblright{}}

\noindent \textbf{Exames de sangue:} Todos dentro da normalidade. Glicemia, ureia, creatinina, sódio, potássio, cálcio, magnésio --- normais. Hormônios tireoidianos --- normais. Sorologias --- todas negativas.

\noindent \textit{\textquotedblleft{}A senhora está saudável\textquotedblright{}\textit{, disse o médico, fechando o prontuário. }\textquotedblleft{}O olho… bem, às vezes acontece. Heterocromia adquirida idiopática. Não tem explicação, mas não mata ninguém.\textquotedblright{}}

\noindent --- \textit{Não tem explicação.}

\noindent Lúcia guardou as palavras. Saiu do hospital carregando a pasta de exames. Atravessou a rua. Sentou num banco de praça. Abriu a pasta. Leu todos os laudos novamente.

\noindent Tudo normal.

\noindent Ela estava saudável.

\noindent Então por que o sonho não parava? Por que o cheiro adocicado aparecia em horários aleatórios? Por que ela acordava toda noite às 3:14?

\noindent \textit{\textquotedblleft{}Não tem explicação.\textquotedblright{}}

\noindent Ela fechou a pasta. Olhou para o céu. O céu estava azul, mas ela não conseguia ignorar a mancha amarela na sua visão periférica.

\noindent --- \textbf{Não tem explicação.}

\noindent Foi aí que ela decidiu: se a ciência não podia responder, ela iria ao lugar onde tudo começou.

\vspace{0.5em}\noindent\rule{\textwidth}{0.5pt}\vspace{0.5em}

\vspace{0.5em}
\noindent\textit{\textbf{↪ Rotas:} \rota{LUC-Entrevista} \rotasep \rota{LUC-05}}